\par\Needspace{6\baselineskip}\medskip
\SourceLeaf{140}\pdfbookmark[2]{100.}{article-100}
\begin{center}100.\end{center}
In solvenda aequatione 13 vel $13^*$ eum casum primo seorsim considerabimus, ubi $z$ obtinet valorem haud magnum, ita ut $Z$ per seriem secundum potestates ipsius $z$ progredientem celeriterque convergentem exprimi possit. Iam fit
\[
\begin{gathered}
(1+2z)\BareRoot{(z+zz)}=z^{\frac12}+\tfrac52z^{\frac32}+\tfrac78z^{\frac52}-\cdots,\\[.5em]
\log(\BareRoot{(1+z)}+\BareRoot z)=z^{\frac12}-\tfrac16z^{\frac32}+\tfrac3{40}z^{\frac52}-\cdots,
\end{gathered}
\]
\noindent adeoque numerator ipsius $Z=\tfrac83z^{\frac32}+\tfrac45z^{\frac52}-\cdots$; denominator autem fit $=2z^{\frac32}+3z^{\frac52}+\cdots$, unde $Z=\tfrac43-\tfrac85z+\cdots$. Ut legem progressionis detegamus, differentiamus aequationem
\[
2(z+zz)^{\frac32}Z=(1+2z)\BareRoot{(z+zz)}-\log(\BareRoot{(1+z)}+\BareRoot z),
\]
\noindent unde prodit omnibus rite reductis
\[
2(z+zz)^{\frac32}\frac{\mathrm dZ}{\mathrm dz}+3Z(1+2z)\BareRoot{(z+zz)}=4\BareRoot{(z+zz)}
\]
\noindent sive
\[
(2z+2zz)\frac{\mathrm dZ}{\mathrm dz}=4-(3+6z)Z,
\]
\noindent unde simili ratione ut in art. 90 deducitur
\[
\begin{aligned}
Z={}&\tfrac43-\frac{4\mathpunct{.}6}{3\mathpunct{.}5}z+\frac{4\mathpunct{.}6\mathpunct{.}8}{3\mathpunct{.}5\mathpunct{.}7}zz\\[.5em]
&-\frac{4\mathpunct{.}6\mathpunct{.}\overset{\mkern2mu .}{8}\mathpunct{.}10}{3\mathpunct{.}5\mathpunct{.}7\mathpunct{.}9}z^3+\frac{4\mathpunct{.}6\mathpunct{.}8\mathpunct{.}10\mathpunct{.}12}{3\mathpunct{.}5\mathpunct{.}7\mathpunct{.}9\mathpunct{.}11}z^4-\text{etc.}
\end{aligned}
\]
\noindent Patet itaque, $Z$ prorsus eodem modo a $-z$ pendere, ut supra in ellipsi $X$ ab $x$: quamobrem si statuimus
\[
Z=\frac1{\tfrac34+\tfrac9{10}(z+\zeta)}\mathpunct{\raisebox{.45ex}{,}}
\]
\noindent determinabitur etiam $\zeta$ perinde per $-z$ ut supra $\xi$ per $x$, ita ut habeatur
\[
\mathrm{14)}\qquad\zeta=\cfrac{\tfrac2{35}zz}{1-\tfrac2{35}z+\cfrac{\tfrac{40}{63}z}{1+\cfrac{\tfrac4{99}z}{1+\cfrac{\tfrac{70}{143}z}{1+\text{etc.}}}}}
\]
\noindent sive
\[
\zeta=\cfrac{\tfrac2{35}zz}{1+\tfrac{18}{35}z+\cfrac{\tfrac4{63}z}{1+\cfrac{\tfrac{40}{99}z}{1+\cfrac{\tfrac{18}{143}z}{1+\text{etc.}}}}}
\]
\noindent\SourceLeaf{141}Hoc modo computati sunt valores ipsius $\zeta$ pro $z=0$ usque ad $z=0{,}3$ per singulas partes millesimas, quos columna tertia tabulae III exhibet.
